\section{Desarrollo matematico extendido}

El aporte teorico que se conserva en anexo es la unificacion de tres lecturas: juego,
mercado y mecanica. El mismo deficit $\delta_k = D_k-S_k$ alimenta payoffs, precios y
requerimientos de wrench. En el modelo reducido, la afirmacion que guia el diseno se
puede formular como una proposicion de alcance limitado.

\paragraph{Proposicion de trabajo.}

\begin{quote}
Si el grafo de comunicacion es conexo a trozos, los payoffs derivan de un potencial
regular, la dinamica Smith es mas rapida que la cinematica de transporte, la formacion
rigida preserva contactos utiles y el campo de navegacion mantiene seguridad, entonces
el sistema acoplado converge practicamente a un conjunto donde las cargas factibles
reciben coaliciones con deficit acotado y residual wrench bajo umbral.
\end{quote}

\paragraph{Hipotesis tecnicas.}
La afirmacion requiere: (i) payoffs localmente Lipschitz, (ii) conectividad suficiente
para que el consenso de ocupacion no acumule sesgo, (iii) separacion de escalas entre
asignacion y movimiento, (iv) saturaciones acotadas, y (v) ausencia de bloqueos
geometricos persistentes. Sin estas condiciones, la proposicion se interpreta como
criterio de diseno, no como teorema cerrado.

\paragraph{Esquema de demostracion.}
La prueba completa queda fuera del cuerpo por alcance. El esquema es:
\begin{enumerate}[leftmargin=2em]
  \item Smith aumenta el potencial hasta un conjunto de Nash del juego poblacional.
  \item El clearing entero perturba la solucion continua con error acotado.
  \item El consenso dinamico mantiene $\xhat$ cerca de la ocupacion real si el grafo no
  pierde conectividad durante demasiado tiempo.
  \item La formacion rigida asigna a cada robot un slot fijo $(r_j,n_j)$ en el marco de
  la carga, por lo que el objetivo local cambia suavemente con la pose del objeto.
  \item El juego vectorial de wrench resuelve una proyeccion acotada sobre
  $\{G_C\lambda: 0\le\lambda\le\bar f\}$ y rechaza coaliciones que no puedan producir
  fuerza y torque.
  \item El campo de navegacion estabiliza la llegada local salvo obstaculos o saturacion.
\end{enumerate}

Esta no es una teoria completamente nueva desde cero: usa dinamicas poblacionales,
consenso, grafos, potenciales, formaciones rigidas y wrench closure. La novedad
defendible esta en ensamblar esas piezas para cargas heterogeneas con una senal
fisico-economica comun y validacion computacional trazable.

\paragraph{Limites de validez.}
La formulacion no demuestra estabilidad global ante colisiones, comunicacion
permanentemente desconectada, contacto fisico real con friccion no modelada ni ejecucion
industrial. Esas extensiones exigen CBF-QP, contacto dinamico o validacion hardware.

\subsection*{Demostraciones minimas por modulo}

\paragraph{M1.1: regla de water-filling.}
En el caso homogeneo, supongase que cada carga activa tiene utilidad marginal decreciente
$p_k(n_k)=v_k-\alpha n_k$ y que los robots pueden redistribuirse entre cargas. En
equilibrio interior de Smith no puede existir un par $(k,\ell)$ con
$p_k(n_k)>p_\ell(n_\ell)$ y robots asignados a $\ell$, porque la dinamica traslada masa
decisional hacia $k$:
\[
  \dot y_{i k} \supset y_{i\ell}[p_k-p_\ell]_+ > 0 .
\]
Por tanto, todas las cargas usadas igualan payoff marginal:
\[
  p_k(n_k)=\lambda,\qquad
  n_k=\frac{v_k-\lambda}{\alpha},
\]
con $\lambda$ escogido para satisfacer $\sum_k n_k=N$. Esta es la lectura tipo
\emph{water-filling} que valida H1.

\paragraph{M1.4: contraejemplo de cardinalidad.}
Sea una carga que requiere torque planar $w^{\rm dem}=[0,0,\tau]^\top$. Dos robots
colocados en la misma cara con normales paralelas pueden cumplir cardinalidad
$|C|=2$, pero su matriz de contactos
\[
  G_C=\begin{bmatrix}
  n_{1x} & n_{2x}\\
  n_{1y} & n_{2y}\\
  r_{1x}n_{1y}-r_{1y}n_{1x} &
  r_{2x}n_{2y}-r_{2y}n_{2x}
  \end{bmatrix}
\]
puede tener rango insuficiente o torque casi nulo. Entonces
\[
  \min_{0\le \lambda\le \bar f}\norm{G_C\lambda-w^{\rm dem}}_2>\epsilon ,
\]
aunque $|C|$ alcance el minimo. Este es el gate G3: quorum numerico no implica capacidad
wrench.

\paragraph{M3: rol fisico como decision adicional.}
Si $h$ indexa slots de contacto, la asignacion completa es
\[
  \max_{z_{ikh}}\sum_{i,k,h} U_{ikh}z_{ikh},
  \qquad
  \sum_{k,h}z_{ikh}\le 1,\qquad
  \sum_i z_{ikh}\le 1 .
\]
El termino $U_{ikh}$ debe mezclar valor de carga, coste de llegada, compatibilidad del
robot con el slot, bateria y reduccion esperada del residual wrench. La version actual
aproxima esta decision en dos pasos: primero elige $k$ con Smith-QR y luego deriva $h$
por ranking de coalicion. Esta aproximacion es ejecutable, pero no equivale a resolver
el problema $z_{ikh}$ completo.

\paragraph{M4--M5: energia y bateria en el potencial.}
El campo unico puede leerse como descenso de una energia compuesta:
\[
  \Phi_i =
  \Phi_i^{\rm tarea}+\Phi_i^{\rm formacion}+\Phi_i^{\rm obstaculo}
  +\Phi_i^{\rm bateria}+\Phi_i^{\rm congestion}.
\]
Si $b_i$ es la bateria normalizada y $\rho_i(q)$ estima coste de retorno, una penalizacion
minima es
\[
  \Phi_i^{\rm bateria} =
  \kappa_b \left[\rho_i(q_i)-b_i+b_{\min}\right]_+^2 .
\]
Asi, el robot deja de ser atractivo para cargas lejanas antes de quedar sin retorno. M5
completo agregaria estaciones $s\in\mathcal S$, colas $Q_s$ y una estrategia adicional de
carga:
\[
  p_{i,s}^{\rm charge}= -\alpha d(q_i,s)-\gamma Q_s+\eta(1-b_i).
\]
Esta memoria usa la primera parte como restriccion de retornabilidad; la planificacion
con estaciones queda como extension.
